\section{The Radical of an Algebra}

Let A be a k-algebra. A left pseudomodule over A is a k-module M endowed with the structure of a left pseudomodule over the pseudoring A (II, Appendix, No.~2, p.~378) such that we have \(a(\lambda x) = \lambda(ax) = (\lambda a)x\) for \(\lambda \in k\) , \(a \in A\) , and \(x \in M\) . We define right pseudomodules over A likewise.

Let M be a left pseudomodule over A; we define a left \(\tilde{A}\) -module structure on M called canonical by setting

\[
(\lambda e + a) x = \lambda x + a x \quad \text { for } \lambda \in k, a \in A, x \in M.
\]

Conversely, every left \(\widetilde{A}\) -module is canonically endowed, by restricting the ring of operators to k and A, with the structure of a left pseudomodule over A. Thus, the species of structures of left pseudomodules over A and of modules over \(\widetilde{A}\) are equivalent (Set Theory, IV, §1, No.~7, p.~262).

Let M be a left pseudomodule over A, and let N be a k-submodule of M. Then N is an \(\widetilde{A}\) -submodule of M if and only if it is stable under the action of A; we then say that N is a sub-pseudomodule of M.

As in the case of rings, we define the left pseudomodule \(A_{s}\) over A and the right pseudomodule \(A_{d}\) . The left (resp. right) ideals of A are the sub-pseudomodules of \(A_{s}\) (resp. \(A_{d}\) ).

Let \(\mathfrak{a}\) be a regular left ideal of A and \(u\) a right unit modulo \(\mathfrak{a}\) . Set \(\mathrm{M} = \mathrm{A}_s / \mathfrak{a}\) , and denote the image of \(u\) in M by \(z\) . We have \(\mathrm{M} = \mathrm{A}z\) , and \(\mathfrak{a}\) is the annihilator of \(z\) .

Conversely, let M be a left pseudomodule over A, and let z be an element of M such that M = Az. Then there exists an element u of A such that z = uz. For every \(a \in A\) , we have \((au - a)z = 0\) , so au - a belongs to the annihilator \(\mathfrak{a}\) of z. Consequently, \(\mathfrak{a}\) is a regular left ideal of A, the element u is a right unit modulo \(\mathfrak{a}\), and when passing to the quotient, the mapping \(a \mapsto az\) defines an isomorphism from \(A_{s}/\mathfrak{a}\) to M.

DEFINITION 2. --- We say that a pseudomodule M over A is simple if we have AM ≠ 0 and if 0 and M are the only sub-pseudomodules of M.

This corresponds to saying that the \(\tilde{A}\) -module M is simple and that its annihilator does not contain A. When A is a ring and M is an A-module, we have AM = M, so Definition 2 coincides with Definition 1 of VIII, p.~45.

PROPOSITION 5. --- a) Let \(\mathfrak{m}\) be a regular maximal left ideal of A. Then the pseudomodule \(A_s / \mathfrak{m}\) is simple, and there exists a nonzero element of \(A_s / \mathfrak{m}\) with annihilator \(\mathfrak{m}\) .

\begin{enumerate}
\def\labelenumi{\alph{enumi})}
\setcounter{enumi}{1}
\item
  Let M be a simple pseudomodule, and let x be a nonzero element of M. We have M = Ax, the annihilator \(\mathfrak{m}\) of x is a regular maximal left ideal of A, and when passing to the quotient, the mapping \(a \mapsto ax\) defines an isomorphism from \(A_{s}/\mathfrak{m}\) to M.
\item
  Let M be a pseudomodule not reduced to 0. Then M is simple if and only if we have M = Ax for every nonzero element x of M.
\end{enumerate}

Let \(\mathfrak{m}\) be a regular maximal left ideal of A. We have seen that there exists a nonzero element \(z\) of \(\mathrm{A}_s / \mathfrak{m}\) with annihilator equal to \(\mathfrak{m}\) and such that \(\mathrm{A}z = \mathrm{A}_s / \mathfrak{m}\) ; in particular, we have \(\mathrm{A}(\mathrm{A}_s / \mathfrak{m}) \neq 0\) . Moreover, every sub-pseudomodule of \(\mathrm{A}_s / \mathfrak{m}\) is of the form \(\mathfrak{n} / \mathfrak{m}\) , where \(\mathfrak{n}\) is a left ideal of A containing \(\mathfrak{m}\) . Since \(\mathfrak{m}\) is maximal, the only possibilities are \(\mathfrak{n} = \mathfrak{m}\) and \(\mathfrak{n} = \mathrm{A}\) , so that the pseudomodule \(\mathrm{A}_s / \mathfrak{a}\) is simple. This proves a).

Under the assumptions of b), the set of elements y of M such that Ay = 0 is a sub-pseudomodule of M different from M and therefore reduced to 0. Consequently, Ax is a nonzero sub-pseudomodule of M, which implies M = Ax. Assertion b) then follows from the remarks before Definition 2.

Let M be a nonzero pseudomodule. Suppose that we have Ax = M for every nonzero element x of M. In particular, we have AM ≠ 0. Let N be a nonzero sub-pseudomodule of M, and let x be a nonzero element of N. We have Ax = M, and therefore N = M. Hence M is simple. Conversely, if M is simple, then we have M = Ax for every \(x \neq 0\) by b).

DEFINITION 3. --- The radical of the k-algebra A, denoted by \(\Re(\mathrm{A})\) , is the intersection of the regular maximal left ideals of A.

When A is a ring, every left ideal of A is regular, so the definition of the radical coincides with Definition 2 of VIII, p.~154.

Examples. --- 1) The radical of the \(k\) -algebra \(\operatorname{End}_k^{\mathrm{f}}(\mathrm{V})\) is reduced to 0 (VIII, p.~436, Example 1). *The same holds for the algebras \(\mathcal{C}_0(\mathrm{T})\) and \(\mathrm{L}^1 (\mathbf{R})\)).*

\begin{enumerate}
\def\labelenumi{\arabic{enumi})}
\setcounter{enumi}{1}
\tightlist
\item
  Let A be a pseudoring in which all elements are nilpotent. The radical of A is equal to A because A is the only regular left ideal.
\end{enumerate}

Proposition 5 immediately implies the following result.

PROPOSITION 6. --- The radical of the algebra A is the intersection of the annihilators of the simple pseudomodules. It is, in particular, a two-sided ideal of A.

PROPOSITION 7. --- The radical of A is the trace on A of the radical of \(\widetilde{A}\) ; it is also equal to the radical of the opposite algebra \(A^{\circ}\) (that is, to the intersection of the regular maximal right ideals of A). If the ring k is without radical, then the radical of A is equal to that of \(\widetilde{A}\) .

The equality \(\Re (\widetilde{\mathrm{A}})\cap \mathrm{A} = \Re (\mathrm{A})\) follows from Proposition 4 of VIII, p.~438, b). Since \(\widetilde{\mathrm{A}}\) and \(\widetilde{\mathrm{A}}^{\circ}\) have the same radical (VIII, p.~156, Corollary 1), the equality \(\Re (\mathrm{A}) = \Re (\mathrm{A}^{\circ})\) follows. If \(k\) is without radical, then the intersection of the maximal left ideals of \(\widetilde{\mathrm{A}}\) containing A is equal to A. Consequently, \(\Re (\widetilde{\mathrm{A}})\) is contained in A and therefore equal to \(\Re (\mathrm{A})\) .

Remark. --- Let x and y be elements of A; we say that x is a left adverse of y, or that y is a right adverse of x, if, in \(\widetilde{A}\) , x - e is a left inverse of y - e, that is, if we have \(x + y = xy\) . By Proposition 7 and Jacobson's theorem (VIII, p.~156, Theorem 1), the radical of A consists of the elements x of A such that ux - e is left invertible in \(\widetilde{A}\) for every u in \(\widetilde{A}\) . Since we have \((a + \lambda e)x - e = ax + \lambda x - e\) (for \(a \in A\) , \(\lambda \in k\) ), the radical of A is the set of elements x of A such that \(ax + \lambda x\) has a left adverse in A for every \(a \in A\) and \(\lambda \in k\) .

